\documentclass[14pt,a4paper,oneside]{extbook}
\usepackage[utf8]{vietnam}
\usepackage{amsmath,amssymb,mathrsfs,fancyhdr,enumerate,multirow,makecell}
\usepackage{tkz-euclide,tikz-3dplot,tikz,tkz-tab}
\usetikzlibrary{arrows.meta,calc,intersections,angles,quotes,backgrounds,shapes.geometric,patterns,through,shadows,positioning}
\usepackage{yhmath}
\usepackage[top=0.7cm, bottom=0.7cm, left=1.0cm, right=1.0cm, headheight=16pt, headsep=3mm, footskip=5mm]{geometry}
\usepackage[hidelinks,unicode]{hyperref}
\usepackage{qrcode}
\usepackage[many]{tcolorbox}
\usepackage{fontawesome}

\pagestyle{fancy}
\fancyhf{}
\fancyhead[L]{\footnotesize\bfseries\color{teal!85!black}\faBook~SỔ TAY KIẾN THỨC TOÁN 4}
\fancyhead[R]{\footnotesize\bfseries\color{blue!85!black}Lớp Toán Thầy Lĩnh -- loptoanthaylinh.vn}
\fancyfoot[C]{\footnotesize\bfseries\color{gray!80!black}-- Sổ tay kiến thức: Bài 68. Ôn tập các phép tính với số tự nhiên --}
\renewcommand{\headrulewidth}{0.4pt}

\begin{document}

% --- BANNER TIÊU ĐỀ BÀI HỌC ---
\noindent
\begin{tcolorbox}[
  enhanced,
  colback=teal!8!white,
  colframe=teal!80!black,
  arc=3mm, boxrule=1pt,
  left=3mm, right=3mm, top=2mm, bottom=2mm,
  halign=center
]
  {\Large\bfseries\color{teal!90!black}\faBookmark\ BÀI 68: ÔN TẬP CÁC PHÉP TÍNH VỚI SỐ TỰ NHIÊN}\\[1.5pt]
  {\footnotesize\itshape\color{black!75} 4 Phép tính, tính chất giao hoán, kết hợp, phân phối và 2 dạng toán lời văn điển hình trên 1 trang A4}
\end{tcolorbox}

\vspace{-2.5mm}

\noindent
\begin{minipage}[t]{0.49\textwidth}\vspace{0pt}
\begin{tcolorbox}[
  enhanced,
  colback=blue!4!white,
  colframe=blue!75!black,
  arc=2mm, boxrule=0.8pt,
  left=2mm, right=2mm, top=1.8mm, bottom=1.8mm,
  title={\footnotesize\bfseries\color{white}\faCogs\ 1. Tính chất các phép tính}
]
{\footnotesize
\begin{itemize}\setlength{\itemsep}{1pt}\setlength{\leftskip}{-2.5mm}
  \item \textbf{Giao hoán:} $a + b = b + a$; \quad $a \times b = b \times a$.
  \item \textbf{Kết hợp:} $(a+b)+c = a+(b+c)$;\\
  $(a \times b) \times c = a \times (b \times c)$.
  \item \textbf{Phân phối nhân với cộng/trừ:}\\
  $a \times (b \pm c) = a \times b \pm a \times c$.
  \item \textbf{Phép chia có dư:} $a = b \times q + r$ ($r < b$).
\end{itemize}
}
\end{tcolorbox}
\end{minipage}
\hfill
\begin{minipage}[t]{0.49\textwidth}\vspace{0pt}
\begin{tcolorbox}[
  enhanced,
  colback=teal!4!white,
  colframe=teal!75!black,
  arc=2mm, boxrule=0.8pt,
  left=2mm, right=2mm, top=1.8mm, bottom=1.8mm,
  title={\footnotesize\bfseries\color{white}\faBook\ 2. Hai dạng toán lời văn kinh điển}
]
{\footnotesize
\begin{itemize}\setlength{\itemsep}{1.5pt}\setlength{\leftskip}{-2.5mm}
  \item \textbf{1. Tìm số trung bình cộng (TBC):}\\
  $\text{TBC} = \text{Tổng các số} : \text{Số các số hạng}$.\\
  $\text{Tổng} = \text{TBC} \times \text{Số các số hạng}$.
  \item \textbf{2. Tìm hai số khi biết Tổng và Hiệu:}\\
  $\text{Số lớn} = (\text{Tổng} + \text{Hiệu}) : 2$.\\
  $\text{Số bé} = (\text{Tổng} - \text{Hiệu}) : 2$.
\end{itemize}
}
\end{tcolorbox}
\end{minipage}

\vspace{-2.5mm}
% --- KHỐI 3: SƠ ĐỒ TƯ DUY TIKZ ---
\noindent
\begin{tcolorbox}[
  enhanced,
  colback=orange!4!white,
  colframe=orange!80!black,
  arc=2mm, boxrule=0.8pt,
  left=2mm, right=2mm, top=1.5mm, bottom=1.5mm,
  title={\footnotesize\bfseries\color{white}\faSitemap\ SƠ ĐỒ TƯ DUY: CÂY CÔNG THỨC 2 DẠNG TOÁN ĐIỂN HÌNH}
]
\centering
\begin{tikzpicture}[
  box/.style={draw, rounded corners=2pt, align=center, font=\scriptsize\bfseries, inner sep=3pt, minimum height=1.05cm},
  arr/.style={-{Stealth[length=2.5mm]}, thick}
]
  \node[box, fill=purple!12, draw=purple!80!black, text width=2.8cm] (root) at (0, 0) {TOÁN LỜI VĂN\\LỚP 4};

  \node[box, fill=blue!10, draw=blue!80!black, text width=4.4cm] (b1) at (5.4, 0.85) {TÌM SỐ TRUNG BÌNH CỘNG\\$\text{TBC} = \text{Tổng} : \text{Số số hạng}$\\$\text{Tổng} = \text{TBC} \times \text{Số số hạng}$};

  \node[box, fill=teal!10, draw=teal!80!black, text width=4.4cm] (b2) at (5.4, -0.85) {TÌM HAI SỐ KHI BIẾT TỔNG \& HIỆU\\$\text{Số lớn} = (\text{Tổng} + \text{Hiệu}) : 2$\\$\text{Số bé} = (\text{Tổng} - \text{Hiệu}) : 2$};

  \draw[arr, blue!80!black] (root.east) -- ++(0.4, 0) |- (b1.west);
  \draw[arr, teal!80!black] (root.east) -- ++(0.4, 0) |- (b2.west);
\end{tikzpicture}
\end{tcolorbox}

\vspace{-2.5mm}
% --- KHỐI 4: VÍ DỤ & MẸO TRÁNH BẪY ---
\noindent
\begin{minipage}[t]{0.49\textwidth}\vspace{0pt}
\begin{tcolorbox}[
  enhanced,
  colback=green!4!white,
  colframe=green!60!black,
  arc=2mm, boxrule=0.8pt,
  left=2mm, right=2mm, top=1.8mm, bottom=1.8mm,
  title={\footnotesize\bfseries\color{white}\faCheckCircle\ Ví dụ thực hành chuẩn mực}
]
{\footnotesize
\textbf{Bài toán (Tổng - Hiệu):}\\
Tổng số tuổi hai anh em là 20 tuổi, anh hơn em 4 tuổi. Tìm tuổi của mỗi người?\\
\textbf{Bài giải:}\\
Tuổi của anh (số lớn) là:\\
$(20 + 4) : 2 = \mathbf{12}$ (tuổi).\\
Tuổi của em (số bé) là:\\
$12 - 4 = \mathbf{8}$ (tuổi).\\
\textbf{Đáp số:} Anh: 12 tuổi; Em: 8 tuổi.
}
\end{tcolorbox}
\end{minipage}
\hfill
\begin{minipage}[t]{0.49\textwidth}\vspace{0pt}
\begin{tcolorbox}[
  enhanced,
  colback=red!4!white,
  colframe=red!70!black,
  arc=2mm, boxrule=0.8pt,
  left=2mm, right=2mm, top=1.8mm, bottom=1.8mm,
  title={\footnotesize\bfseries\color{white}\faExclamationTriangle\ Mẹo tránh bẫy trắc nghiệm}
]
{\footnotesize
\begin{itemize}\setlength{\itemsep}{1.5pt}\setlength{\leftskip}{-2.5mm}
  \item \textbf{Bẫy thiếu dấu ngoặc ở phép chia 2:}\\
  Viết $20 + 4 : 2 = 22$ là SAI! Bắt buộc phải đóng mở ngoặc: $(20 + 4) : 2 = 12$.
  \item \textbf{Bẫy điều kiện số dư:}\\
  Trong phép chia có dư, số dư lớn nhất luôn bằng $\text{Số chia} - 1$. Nếu dư $\ge$ chia $\to$ Phép chia thực hiện sai!
\end{itemize}
}
\end{tcolorbox}
\end{minipage}

\end{document}
